\documentclass[aspectratio=169]{beamer}
\usepackage[T1]{fontenc}
\usepackage[utf8]{inputenc}
\usepackage{hyperref}
\usetheme{Madrid}
\setbeamertemplate{navigation symbols}{}

\title{GENRTL Progress Update}
\subtitle{SPEC-to-RTL Workflow Comparison}
\author{Jason Pan}
\date{\today}

\begin{document}

\begin{frame}
  \titlepage
\end{frame}

\begin{frame}{Project Goal}
  \begin{itemize}
    \item Compare two Cursor-Agent flows for hardware design:
      \begin{itemize}
        \item Workflow A: SPEC $\rightarrow$ Verilog RTL
        \item Workflow B: SPEC $\rightarrow$ refine SPEC(agent1)$\rightarrow$ timing plan(agent2) $\rightarrow$ RTL(agent3)
        \begin{figure}
            \centering
            \includegraphics[width=0.6\linewidth]{agent2_flowChart.png}
            \label{fig:placeholder}
        \end{figure}
      \end{itemize}
    \item Benchmark/evaluation target: RTLLM + VCS simulation
    \item Fairness setting: Pass@1 (single generation, no testbench/log feedback during generation)
  \end{itemize}
\end{frame}

\begin{frame}{Workflow B Roles (from SKILL.en.md)}
\small

\begin{columns}[t] % [t] 表示頂端對齊，可讓兩邊文字對齊更美觀
  
  % 左半邊：Agent1
  \begin{column}{0.48\textwidth}
    \textbf{Agent1 — Refine SPEC}
    \vspace{0.2cm} % 稍微留點空隙
    
    \centering
    \includegraphics[width=1\linewidth]{agent1.png}
  \end{column}

  % 右半邊：Agent2
  \begin{column}{0.48\textwidth}
    \textbf{Agent2 — Timing \& Structure}
    \vspace{0.2cm}
    
    \centering
    \includegraphics[width=1\linewidth]{agent2.png}
  \end{column}

\end{columns}

\end{frame}

\begin{frame}{What Has Been Built}
  \begin{itemize}
    \item Project skeleton and directory conventions completed
    \item Cursor setup completed:
      \begin{itemize}
        \item Skills: \texttt{@rtl-workflow-a}, \texttt{@rtl-pipeline-workflow-b}
        \item Rules for path constraints and Pass@1 fairness
        \item \texttt{.cursorignore} to avoid indexing reference/golden RTL during generation
      \end{itemize}
    \item Docs updated: README/HANDOFF/AGENTS/prompts/spec/experiments templates
  \end{itemize}
\end{frame}



\begin{frame}[fragile]{Repository Structure}
\begin{figure}
    \centering
    \includegraphics[width=0.7\linewidth]{tree.png}
    \label{fig:placeholder}
\end{figure}
\end{frame}


\begin{frame}{Current Experiment Status}
  \begin{itemize}
    \item Test case tried: \texttt{adder\_pipe\_64bit}
    \item Workflow A and Workflow B both generated RTL outputs
    \item VCS compile stage can complete
    \item Functional simulation currently pass Workflow B only.
  \end{itemize}
\end{frame}


\begin{frame}{Next Steps}
  \begin{itemize}
    \item Run A/B across multiple RTLLM problems with same protocol
    \item Record all runs in \texttt{experiments/}:
      \begin{itemize}
        \item model, skill used, SPEC path, VCS compile/sim result
      \end{itemize}
    \item Add batch evaluation script (\texttt{scripts/run\_vcs.sh}) and benchmark manifest
    \item Compare WorkflowA and WorkflowB with other models in paper.
  \end{itemize}
\end{frame}

\begin{frame}{References}
\footnotesize
\begin{thebibliography}{9}
\bibitem{rtllm}
Yao Lu, Shang Liu, Qijun Zhang, and Zhiyao Xie, \\
\emph{RTLLM: An Open-Source Benchmark for Design RTL Generation with Large Language Model}, ASP-DAC 2024. \\
\url{https://github.com/hkust-zhiyao/RTLLM}

\bibitem{openllmrtl}
Shang Liu, Yao Lu, Wenji Fang, Mengming Li, and Zhiyao Xie, \\
\emph{OpenLLM-RTL: Open Dataset and Benchmark for LLM-Aided Design RTL Generation}, ICCAD 2024 (Invited).

\bibitem{genrtlrepo}
GENRTL Project Repository, \\
\url{https://github.com/jasonpan0930/GENRTL}
\end{thebibliography}
\end{frame}


\end{document}

